% answers_short.tex —— 2006 年·普通高等学校招生考试·湖北卷理科 · 简短答案速查
% 由 tools/make_short_answers.py 自 answers.tex 的 \gksolhead/\gkanswer 提取，
% 并与 2006/2006_数学真题_湖北理_答案要点.md 逐题比对后生成，勿手改（改源 MD 后重跑）。

\begin{gknotice}
\textbf{使用说明}：本节只给最终答案便于快速核对，完整推导见解析版。答案为编写者独立推导，\textbf{非官方标准答案}。
\end{gknotice}

\gkpart{选择题}
\begin{center}\begin{tabular}{ccccccccccc}
\toprule
题号 & 1 & 2 & 3 & 4 & 5 & 6 & 7 & 8 & 9 & 10 \\
\midrule
答案 & B & D & A & B & C & D & D & B & C & A \\
\bottomrule
\end{tabular}\end{center}

\gkpart{填空题}
\begin{itemize}[leftmargin=2.6em]
  \item[\textbf{11.}] $4$。
  \item[\textbf{12.}] $0.94$。
  \item[\textbf{13.}] $8$ 或 $-18$。
  \item[\textbf{14.}] $20$。
  \item[\textbf{15.}] $1/2$。
  \item[\textbf{16.}] (1) 最大值 $2+\sqrt2$，最小正周期 $\pi$；(2) 最短平移向量为 $(-\pi/8,-2)$
\end{itemize}

\gkpart{解答题}
\begin{itemize}[leftmargin=2.6em]
  \item[\textbf{17.}] (1) $a_n=6n-5$；(2) U/PENDING；
  \item[\textbf{18.}] (1) $m=1/3$；(2) 存在；
  \item[\textbf{19.}] (1) 约 $526$ 人；(2) 约 $83.1$ 分
  \item[\textbf{20.}] (1) $\frac{x^2}{4}+\frac{y^2}{3}=1$；(2) 见解析版
  \item[\textbf{21.}] (1) $b=-2a-3$；(2) $0<a<3/2$
\end{itemize}
